
\subsection{滤过 $A_{n,K}$ 同伦的构造}
\label{subsec:7-2-10}
\liftIndex{I258}%

本小节旨在证明命题 4.6.37 与定理 4.6.25 的
$A_{n,K}$ 版本。定理 4.6.25 的
$A_\infty$ 版本将在第 7.2.13 小节完成证明。
这里沿用第 4.6.2 小节的记号。

先固定本小节的记号。固定 $\lambda$，
如命题~\ref{prop:7-2-100} 的证明那样，写
$C_1(L;\Lambda_{0,\mathrm{nov}})=C_1(L,\lambda)$，
并写
$C_2(L';\Lambda_{0,\mathrm{nov}})=C_2(L',\lambda')$。
\translatornote{原书在这处等式左、右两边均将 $L'$ 误排为 $L$；
由紧接的生成元及后文目标空间改正。}
可数链集
$\mathcal X_1(L)=\mathcal X_1(L,\lambda)$、
$\mathcal X_2(L')=\mathcal X_2(L',\lambda')$
分别生成它们。下文不再涉及
$\mathcal X_1(L)$、$\mathcal X_2(L')$
中元素的代。

使用 (4.6.26) 中的双参数殆复结构族
$\{J_{\rho,s}\}_{\rho,s}$，并置
\[
\{J_\rho^{s=0}\}_\rho=\{J_{\rho,0}\}_\rho,
\qquad
\{J_\rho^{s=1}\}_\rho=\{J_{\rho,1}\}_\rho.
\]
取第 4.6.2 小节中的辛微分同胚族
$\psi_s\colon(M,L)\to(M',L')$，满足
$\psi_0=\psi$、$\psi_1=\psi'$。
用 $\{J_\rho^{s=0}\}_\rho$ 与
$\psi=\psi_0$ 进行第~\ref{subsec:7-2-9} 小节的构造，
得到 $\mathcal X^{s=0}_{3,g}(L')$（$g\le\mathcal K$）
以及
\[
\mathfrak s'_{s=0,\mathfrak d,k,\beta,\vec P'},
\qquad
\mathfrak s^{\operatorname{top}(\rho)}_{s=0,k,\beta,\vec P}.
\]
这里也省略记号中的 $\lambda,\lambda'$。
它们在由 $\mathcal X^{s=0}_{3,(g)}(L')$ 生成的
$C^{s=0}_3(L';\Lambda_{0,\mathrm{nov}})$ 上
诱导滤过 $A_{n,K}$ 代数结构
$\mathfrak m^{J',\mathfrak s''_{s=0}}$，
并给出滤过 $A_{n,K}$ 同态
\begin{equation}
\mathfrak f_{s=0}^{\psi_0,\{J^{s=0}_\rho\}_\rho,
\mathfrak s^{\operatorname{top}(\rho)}_{s=0}}
\colon
(C_1(L;\Lambda_{0,\mathrm{nov}}),\mathfrak m^{J,\mathfrak s}_k)
\longrightarrow
(C^{s=0}_3(L';\Lambda_{0,\mathrm{nov}}),
\mathfrak m^{J',\mathfrak s''_{s=0}}_k).
\tag{7.2.133}\label{eq:7-2-133}
\end{equation}
这里 $(n,K)=(n(\lambda),K(\lambda))$，
二者在 $\lambda\to\infty$ 时趋于无穷。
严格说，$C^{s=0}_3(L';\Lambda_{0,\mathrm{nov}})$
应写作 $C^{s=0}_{3,(g)}(L';\Lambda_{0,\mathrm{nov}})$；
如同省略 $\lambda$，我们也为简便省略 $(g)$。
还须先构造 $\mathfrak f^{\mathrm{geo}}$，
再像构造 $\mathfrak m^{\mathrm{geo}}$ 时那样修正它；
这里不再讨论细节。

接着用 $\{J^{s=1}_\rho\}_\rho$ 与
$\psi'=\psi_1$ 进行第~\ref{subsec:7-2-9} 小节的构造，
得到 $\mathcal X^{s=1}_{3,g}(L')$
（$g\le\mathcal K=\mathcal K_\lambda$）以及
\[
\mathfrak s''_{s=1,\mathfrak d,k,\beta,\vec P'},
\qquad
\mathfrak s^{\operatorname{top}(\rho)}_{s=1,k,\beta,\vec P}.
\]
它们在由 $\mathcal X^{s=1}_{3,(g)}(L')$ 生成的
$C^{s=1}_3(L';\Lambda_{0,\mathrm{nov}})$ 上
定义滤过 $A_{n,K}$ 结构
$\mathfrak m^{J',\mathfrak s''_{s=1}}$，
并给出滤过 $A_{n,K}$ 同态
\begin{equation}
\mathfrak f_{s=1}^{\psi_1,\{J^{s=1}_\rho\}_\rho,
\mathfrak s^{\operatorname{top}(\rho)}_{s=1}}
\colon
(C_1(L;\Lambda_{0,\mathrm{nov}}),\mathfrak m^{J,\mathfrak s}_k)
\longrightarrow
(C^{s=1}_3(L';\Lambda_{0,\mathrm{nov}}),
\mathfrak m^{J',\mathfrak s''_{s=1}}_k).
\tag{7.2.134}\label{eq:7-2-134}
\end{equation}
由 (7.2.133)、(7.2.134) 分别同下列嵌入
的同伦逆复合：
\[
\begin{aligned}
C_2(L';\Lambda_{0,\mathrm{nov}})
&\longrightarrow C^{s=0}_3(L';\Lambda_{0,\mathrm{nov}}),\\
C_2(L';\Lambda_{0,\mathrm{nov}})
&\longrightarrow C^{s=1}_3(L';\Lambda_{0,\mathrm{nov}}),
\end{aligned}
\]
可得
\[
\mathfrak g_{s=0}^{\psi_0,\{J^{s=0}_\rho\}_\rho,
\mathfrak s^{\operatorname{top}(\rho)}_{s=0}}
\quad\text{与}\quad
\mathfrak g_{s=1}^{\psi_1,\{J^{s=1}_\rho\}_\rho,
\mathfrak s^{\operatorname{top}(\rho)}_{s=1}}.
\]
我们要构造二者之间的 $A_{n,K}$ 同伦。
由于 (7.2.133) 与 (7.2.134) 的目标不同，
还须构造在二者之间插值的滤过 $A_{n,K}$ 代数
$(C([0,1]\times L')\widehat\otimes
\Lambda_{0,\mathrm{nov}},\mathfrak m)$。
下文把 $\mathcal X^{s=0}_{3,g}(L')$、
$\mathcal X^{s=1}_{3,g}(L')$ 简写为
$\mathcal X^{s=0}_3(L')$、$\mathcal X^{s=1}_3(L')$，
省去表示元素代数层级的下标 $g$。
本小节不再增加从 $L$ 或 $L'$ 选取的链数，
而是固定链集；因此，多值截面
$\mathfrak s'_{s=0,\mathfrak d,k,\beta,\vec P'}$、
$\mathfrak s''_{s=1,\mathfrak d,k,\beta,\vec P'}$
中的 $\mathfrak d$ 也省去，分别写作
$\mathfrak s'_{s=0,k,\beta,\vec P'}$、
$\mathfrak s''_{s=1,k,\beta,\vec P'}$。
定义这些多值截面的模空间记号中也省去 $\mathfrak d$。

回忆第 4.6 节对 $[0,1]\times L'$ 上的奇异单形
引入的条件如下。
\begin{statement}{定义}{4.6.30（重述）}
设 $\mathcal P=(|\mathcal P|,f)$ 是光滑奇异单形，
$f\colon|\mathcal P|\to[0,1]\times L'$。
若对 $s_0=0,1$ 都满足下列条件，
则称 $\mathcal P$ 是 $[0,1]\times L'$ 上的
\emph{适配奇异单形}：\liftIndex{I259}%
\translatornote{原书这处称“$[0,1]\times L$ 上的适配单形”，
但映射靶及以下各条件均为 $[0,1]\times L'$，故补撇号。}
\begin{enumerate}[label=\textup{(4.6.31.\arabic*)},leftmargin=5.9em]
 \item $f^{-1}(\{s_0\}\times L')\cap|\mathcal P|$
 或等于 $|\mathcal P|$，或为空，
 或仅由 $|\mathcal P|$ 的一个面（余维任意）组成。
 \item 若 $f^{-1}(\{0\}\times L')\cap|\mathcal P|$
 非空，则
 $|\mathcal P|\subset[0,1/3)\times L'$；
 若 $f^{-1}(\{1\}\times L')\cap|\mathcal P|$
 非空，则
 $|\mathcal P|\subset(2/3,1]\times L'$。
 \item 设 $\Delta_a^m$ 是 $|\mathcal P|$ 的一个面，
 $\Delta_b^{m'}$ 是包含它的另一个面。
 若 $x\in\Delta_a^m\subset
 f^{-1}(\{s_0\}\times L')$，
 而 $\Delta_b^{m'}$ 不包含于
 $f^{-1}(\{s_0\}\times L')$，
 则任意
 $N\in T_x\Delta_b^{m'}\setminus
 T_x(\partial\Delta_b^{m'})$
 在 $f_*(N)$ 下的 $[0,1]$ 分量非零。
\end{enumerate}
\end{statement}

把 $m$ 维奇异链视为次数
$d=1+\dim L'-m$ 的上链。
记 $S_+^d([0,1]\times L')$ 为次数 $d$
的上链组成的阿贝尔群。
令 $S_+([0,1]\times L')$ 为
$[0,1]\times L'$ 上适配奇异单形自由生成的
阿贝尔群；容易看出它是上链复形。
存在典范包含
\[
S(\{0,1\}\times L')\subset S_+([0,1]\times L'),
\]
其中 $S(\{0,1\}\times L')$ 是光滑奇异链复形，
也是 $S_+([0,1]\times L')$ 的子复形。
定义商复形
\[
S([0,1]\times L')
=\frac{S_+([0,1]\times L')}
{S(\{0,1\}\times L')}.
\]
故它以不包含于 $\{0,1\}\times L'$
的适配奇异单形为自由基。

现在定义评价映射
\[
\operatorname{Eval}_{s=s_0}\colon
S([0,1]\times L')\longrightarrow S(L').
\]
若 $f^{-1}(\{s_0\}\times L')\cap|\mathcal P|$
在 $|\mathcal P|$ 中不是余维一，令
$\operatorname{Eval}_{s=s_0}(\mathcal P)=0$。
若它是余维一，令
\[
\operatorname{Eval}_{s=s_0}(\mathcal P)
=(-1)^{s_0+c+1}
(\Delta_c^{m-1},f_c\circ j_c),
\]
其中
$\Delta_c^{m-1}=(v_0,\ldots,\widehat v_c,\ldots,v_m)$
就是该面，
$f_c=f|_{\Delta_c^{m-1}}$，而
$j_c\colon\Delta^{m-1}\to\Delta_c^{m-1}$
是由下式给出的单纯映射：
\[
j_c(v_i)=
\begin{cases}
v_i,&i<c,\\
v_{i+1},&i\ge c.
\end{cases}
\]
\translatornote{原书此处第二分支写 $i>c$，
于是 $j_c(v_c)$ 无定义；按删去第 $c$ 个顶点的标准面嵌入改为 $i\ge c$。}
这里使用以链的余维定义次数的上同调记号，
所以 $\operatorname{Eval}_{s=s_0}$ 的次数为零。

下面把 $[0,1]$ 分成
$[0,1/4]$、$[1/4,3/4]$ 与 $[3/4,1]$，
用棱柱型单纯剖分\liftIndex{I260}%
定义同态
$\operatorname{Incl}_{1,\beta_0}\colon
S(L')\to S([0,1]\times L')$。
令 $[s,s']$ 是上述三个区间之一。
对 $\Delta^m=(v_0,\ldots,v_m)$，
取标准的棱柱单纯剖分
\[
[s,s']\times\Delta^m
=\coprod_{i=0}^{m}
(\underline v_0,\ldots,\underline v_i,
\overline v_i,\ldots,\overline v_m),
\]
其中 $\underline v_i=(s,v_i)$、
$\overline v_i=(s',v_i)$。
定义单纯映射
\[
j_i\colon\Delta^{m+1}=(v_0,\ldots,v_{m+1})
\longrightarrow
(\underline v_0,\ldots,\underline v_i,
\overline v_i,\ldots,\overline v_m)
\]
为
\[
j_i(v_j)=
\begin{cases}
\underline v_j,&j\le i,\\
\overline v_{j-1},&j>i.
\end{cases}
\]
于是对奇异单形 $\sigma\colon\Delta^m\to L'$，定义
\begin{equation}
\begin{aligned}
\operatorname{Incl}_{1,\beta_0}(\sigma)
={}&\sum_{i=0}^{m}(-1)^i
(\operatorname{id}_{[0,1/4]}\times\sigma)\circ j_i\\
&+\sum_{i=0}^{m}(-1)^i
(\operatorname{id}_{[1/4,3/4]}\times\sigma)\circ j_i\\
&+\sum_{i=0}^{m}(-1)^i
(\operatorname{id}_{[3/4,1]}\times\sigma)\circ j_i.
\end{aligned}
\tag{7.2.135}\label{eq:7-2-135}
\end{equation}
按上述分段方式，特别是 $1/4<1/3$，
棱柱剖分中每个单形在
$\operatorname{Incl}_{1,\beta_0}$ 下的像
都是适配单形。

回忆 $C(L';\Lambda_{0,\mathrm{nov}})$ 上的
$\mathfrak m_{1,\beta_0}=(-1)^n\partial$，
$n=\dim L'$。
在 $S([0,1]\times L';\Lambda_{0,\mathrm{nov}})$ 上
置 $\mathfrak m_{1,\beta_0}=(-1)^{n+1}\partial$
（模 $S(\{0,1\}\times L';\Lambda_{0,\mathrm{nov}})$）。
注意 $n+1=\dim([0,1]\times L')$，故易见
\[
\mathfrak m_{1,\beta_0}\circ
\operatorname{Incl}_{1,\beta_0}
=\operatorname{Incl}_{1,\beta_0}\circ
\mathfrak m_{1,\beta_0}.
\]
滤过 $A_\infty$ 代数结构的符号相容性
见第 8.9.2 小节。

\begin{lemma}{7.2.136}\label{lem:7-2-136}
\noindent\textup{(7.2.137.1)}\quad
$\operatorname{Eval}_{s=s_0}$ 对
$\mathfrak m_{1,\beta_0}$ 是链映射。

\noindent\textup{(7.2.137.2)}\quad
$\operatorname{Eval}_{s=0}\oplus
\operatorname{Eval}_{s=1}\colon
S([0,1]\times L')\to S(L')\oplus S(L')$
是满射。

\noindent\textup{(7.2.137.3)}\quad
$\operatorname{Eval}_{s=s_0}$ 与
$\operatorname{Incl}_{1,\beta_0}$
都是链同伦等价。
\end{lemma}

\begin{proof}
先证明 (7.2.137.1)。
令 $(\mathcal P,f)$ 是不包含于
$\{0,1\}\times L'$ 的 $m$ 维适配奇异单形。
分四种情况讨论：
\begin{align*}
\textup{(7.2.138.1)}\quad&
\dim\bigl(f^{-1}(\{0\}\times L')\cap|\mathcal P|\bigr)=m,\\
\textup{(7.2.138.2)}\quad&
\dim\bigl(f^{-1}(\{0\}\times L')\cap|\mathcal P|\bigr)<m-2,\\
\textup{(7.2.138.3)}\quad&
\dim\bigl(f^{-1}(\{0\}\times L')\cap|\mathcal P|\bigr)=m-1,\\
\textup{(7.2.138.4)}\quad&
\dim\bigl(f^{-1}(\{0\}\times L')\cap|\mathcal P|\bigr)=m-2.
\end{align*}
逐种证明
$\partial(\operatorname{Eval}_{s=s_0}(\mathcal P))
=-\operatorname{Eval}_{s=s_0}(\partial\mathcal P)$；
这等价于
\[
\mathfrak m_{1,\beta_0}\circ
\operatorname{Eval}_{s=s_0}(\mathcal P)
=\operatorname{Eval}_{s=s_0}\circ
\mathfrak m_{1,\beta_0}(\mathcal P).
\]
对 (7.2.138.1)、(7.2.138.2)，两端均为零。
对 (7.2.138.3)，由 (4.6.31.1)，
$f^{-1}(\{0\}\times L')=\Delta_c^{m-1}$，
其中 $\Delta_c^{m-1}$ 是 $|\mathcal P|$ 的余维一面。
于是对某个 $c=0,\ldots,m$，有
\[
\operatorname{Eval}_{s=s_0}(\mathcal P)
=(-1)^{s_0+c+1}
(\Delta_c^{m-1},f_c\circ j_c),
\qquad f_c:=f|_{\Delta_c^{m-1}}.
\]
由适配奇异单形的条件，
其余每个面
$(\Delta_b^{m-1},f|_{\Delta_b^{m-1}})$
（$b\ne c$）均不包含于 $\{0\}\times L'$。
记该面为 $\partial_b\mathcal P$，则
$\partial\mathcal P=
\sum_{b=0}^{m}(-1)^b\partial_b\mathcal P$，从而
\[
\begin{aligned}
\operatorname{Eval}_{s=s_0}(\partial\mathcal P)
&=\operatorname{Eval}_{s=s_0}
\left(\sum_{b=0}^{m}(-1)^b\partial_b\mathcal P\right)\\
&=\operatorname{Eval}_{s=s_0}
\bigl((-1)^c\partial_c\mathcal P\bigr)
+\operatorname{Eval}_{s=s_0}
\left(\sum_{b\ne c}(-1)^b\partial_b\mathcal P\right)\\
&=\sum_{b\ne c}(-1)^b
\operatorname{Eval}_{s=s_0}(\partial_b\mathcal P).
\end{aligned}
\]
第一项为零，因为按 $S_+([0,1]\times L')$
的定义，$\partial_c\mathcal P\subset\{s_0\}\times L'$。
另一方面，按单形面的边界符号，有
\[
\begin{aligned}
\operatorname{Eval}_{s=s_0}(\partial\mathcal P)
&=(-1)^{s_0+1}
\Bigl\{\sum_{b<c}(-1)^{b+c}
(\partial_c\partial_b\Delta^m,f)\\
&\hspace{3em}+\sum_{b>c}(-1)^{b+c-1}
(\partial_{c-1}\partial_b\Delta^m,f)\Bigr\}\\
&=(-1)^{s_0+1}
\Bigl\{\sum_{b<c}(-1)^{b+c}
(\partial_{b-1}\partial_c\Delta^m,f)\\
&\hspace{3em}+\sum_{b>c}(-1)^{b+c-1}
(\partial_b\partial_{c-1}\Delta^m,f)\Bigr\}\\
&=(-1)^{s_0+1}
\sum_{b\ne c}(-1)^{b+1}(-1)^c
(\partial_b\partial_c\Delta^m,f)\\
&=-\sum_{b\ne c}(-1)^b
\partial_b\operatorname{Eval}_{s=s_0}(\mathcal P)\\
&=-\partial\bigl(
\operatorname{Eval}_{s=s_0}(\mathcal P)\bigr).
\end{aligned}
\]
这就得到所需等式。
对 (7.2.138.4)，存在两个余维一面
$\Delta_i^{m-1},\Delta_j^{m-1}$，使
\[
f^{-1}(\{0\}\times L')\cap\partial\Delta_i^{m-1}
=-f^{-1}(\{0\}\times L')\cap\partial\Delta_j^{m-1}
=f^{-1}(\{0\}\times L')\cap|\mathcal P|.
\]
因此两端同样为零，(7.2.137.1) 得证。

(7.2.137.2) 显然，略去证明。
现在证明 (7.2.137.3)。
由于
\[
\operatorname{Eval}_{s=s_0}\circ
\operatorname{Incl}_{1,\beta_0}=\mathrm{id},
\]
$\operatorname{Eval}_{s=s_0}$ 在同调上诱导满射。
只需证明其在同调上也是单射。
令
\[
\mathcal C=\sum_{i\in I}a_i\mathcal P_i
\in S^d([0,1]\times L')
\]
是一个闭链，且
$\operatorname{Eval}_{s=0}(\mathcal C)=\partial\mathcal D$
对某条链 $\mathcal D\in S(L')$ 成立。
把 $\mathcal P_i$ 表示为
$S_+^d([0,1]\times L')$ 中的适配奇异单形，
并经由明显的包含
$S(L')\to S(\{0\}\times L')
\subset S_+^d([0,1]\times L')$
把 $\mathcal D$ 看作 $S_+^d([0,1]\times L')$ 的元素。
于是 $\mathcal C-\mathcal D$ 满足
\[
\partial(\mathcal C-\mathcal D)
\in S(\{1\}\times L').
\]
相对同调群
$H([0,1]\times L';\{1\}\times L')$ 为零，
故存在一条奇异链 $\mathcal E$，
使其边界为 $\mathcal C-\mathcal D$。
但 $\mathcal E$ 未必属于
$S_+^{d-1}([0,1]\times L')$。
为找到其中的链，作如下处理。
设 $I_0\subset I$ 由满足
$\mathcal P_i\cap(\{0\}\times L')\ne\varnothing$
的指标 $i$ 组成，并写
$\mathcal P_i=(|\mathcal P_i|,f_i)$。
由其适配性，尤其是 (4.6.31.2)，
存在 $\epsilon>0$，使
\[
f_i(|\mathcal P_i|)
\subset[0,1/3-\epsilon)\times L'
\qquad(i\in I_0).
\]
写 $f_i=(f_{i,s},f_{i,L'})$，定义
\begin{equation}
\widehat f_i(x,t)=
\left(\frac{\epsilon t}{2}+f_{i,s}(x),
f_{i,L'}(x)\right).
\tag{7.2.139}\label{eq:7-2-139}
\end{equation}
用此映射及 $|\mathcal P_i|\times[0,1]$
的棱柱剖分，得到奇异链
$\Psi(\mathcal P_i)\in S_+([0,1]\times L')$。
容易看出
$\partial\Psi(\mathcal P_i)-
\Psi(\partial\mathcal P_i)-\mathcal P_i$
的支集不与 $\{0\}\times L'$ 相交。
类似定义 $\Psi(\mathcal D)$，并令
$\mathcal E_1=
\sum_{i\in I_0}a_i\Psi(\mathcal P_i)-
\Psi(\mathcal D)$。
于是 $\mathcal C-\mathcal D-\partial\mathcal E_1$
不与 $\{0\}\times L'$ 相交，且满足
\[
\partial(\mathcal C-\mathcal D-
\partial\mathcal E_1)
=\partial(\mathcal C-\mathcal D)
\in S(\{1\}\times L').
\]
现在容易找到 $\mathcal E_2$，使
\[
\mathcal C-\mathcal D-
\partial\mathcal E_1-
\partial\mathcal E_2
\in S(\{1\}\times L'),
\]
并使 $\mathcal E=\mathcal E_1+\mathcal E_2$
属于 $S_+^{d-1}([0,1]\times L')$。
这就证明了 (7.2.137.3)。
\end{proof}

现在引入 $[0,1]\times L'$ 上的可数奇异单形集族
$\mathcal X_g([0,1]\times L')$（$g=1,2,\ldots$）；
其条件将在下面的条件 7.2.141 给出。
为此还需一个定义。

\begin{definition}{7.2.140}\label{def:7-2-140}
对 $P\in\mathcal X_1(L)$，定义
$[0,1]\times L'$ 上的链 $\{\psi_s\}_{s*}(P)$ 如下。
设 $P=(|P|,f)$，其中把 $|P|$ 看作带角流形，
$f\colon|P|\to L$ 光滑。
定义
$\{\psi_s\}_s\circ f\colon
[0,1]\times|P|\to[0,1]\times L'$ 为
\[
(\{\psi_s\}_s\circ f)(s,x)
=(s,\psi_s(f(x))).
\]
对 $[0,1]\times|P|$ 取棱柱剖分，
其中仍将 $[0,1]$ 分成
$[0,1/4]$、$[1/4,3/4]$、$[3/4,1]$，
便得到 $[0,1]\times L'$ 上的奇异链，
记作 $\{\psi_s\}_{s*}(P)$。
\end{definition}

在下一条件中，把 $\mathcal X^{s=s_0}_3(L')$
的元素同 $S(\{s_0\}\times L')$ 中的元素等同，
从而可视为 $S_+([0,1]\times L')$ 的子集。

\begin{statement}{条件}{7.2.141}\label{cond:7-2-141}
记
\[
\mathcal X_{(g_0)}([0,1]\times L')
=\bigcup_{g\le g_0}\mathcal X_g([0,1]\times L'),
\]
并以 $C_{(g_0)}([0,1]\times L';\mathbb Q)$
表示由它张成的 $\mathbb Q$ 向量空间。

\noindent\textup{(7.2.142.0)}\quad
$\mathcal X_{00}([0,1]\times L')$
是下列单纯剖分中出现、且不包含于
$\{0,1\}\times L'$ 的奇异单形的并：
对 $P'\in\mathcal X_2(L')$ 的
$[0,1]\times P'$ 的剖分，以及对
$P\in\mathcal X_1(L)$ 的
$\{\psi_s\}_{s*}(P)$ 的剖分。
定义 $\operatorname{Incl}_{1,\beta_0}$
与定义~\ref{def:7-2-140} 时已固定这些剖分。

\noindent\textup{(7.2.142.1)}\quad
$\mathcal X_g([0,1]\times L')$
是 $[0,1]\times L'$ 上光滑奇异单形的可数集。

\noindent\textup{(7.2.142.2)}\quad
包含
$C_{(g)}([0,1]\times L';\mathbb Q)
\to S([0,1]\times L';\mathbb Q)$
诱导上同调同构。

\noindent\textup{(7.2.142.3)}\quad
$\mathcal X_g([0,1]\times L')$ 的每个元素的每个面，
均属于
\[
\mathcal X_{(g)}([0,1]\times L')
\cup\mathcal X_3^{s=0}(L')
\cup\mathcal X_3^{s=1}(L').
\]

\noindent\textup{(7.2.142.4)}\quad
$\mathcal X_g([0,1]\times L')$ 的元素
满足条件~\ref{def:4-6-30}。

\noindent\textup{(7.2.142.5)}\quad
$\mathcal X_g([0,1]\times L')$ 中元素的
$\operatorname{Eval}_{s=s_0}$ 像
属于 $\mathcal X^{s=s_0}_3(L')$。
\end{statement}

(7.2.142.1)--(7.2.142.3) 是性质 7.2.27 的对应版本。
称 $\mathcal X_g([0,1]\times L')$ 的元素具有
\emph{代 $g$}。
以后有时把
$\operatorname{Eval}_{s=s_0}(\mathcal P)$
简记为 $\mathcal P|_{s=s_0}$。

现在回顾将用到的两类模空间。
第一类是 (4.6.29) 定义的
\[
\mathcal M_{k+1}^{\mathrm{main}}
(M',L',\{J_{\rho,s}\}_{\rho,s};
\beta;\operatorname{top}(\rho),
\operatorname{twp}(s);\vec P),
\]
其中 $\vec P=(P_1,\ldots,P_k)$，
$P_i\in\mathcal X_1(L)$。
为简化记号，写作
\begin{equation}
\mathcal M_{k+1}^{\mathrm{main}}
(\{J_{\rho,s}\}_{\rho,s};\beta;
\operatorname{top}(\rho),\operatorname{twp}(s);\vec P).
\tag{7.2.143}\label{eq:7-2-143}
\end{equation}
其评价映射为
\begin{equation}
ev_0^+=(ev_s,ev_0)\colon
\mathcal M_{k+1}^{\mathrm{main}}
(\{J_{\rho,s}\}_{\rho,s};\beta;
\operatorname{top}(\rho),\operatorname{twp}(s);\vec P)
\longrightarrow[0,1]\times L'.
\tag{7.2.144}\label{eq:7-2-144}
\end{equation}
此模空间的仓西结构可选得使
\[
ev\colon
\mathcal M_{k+1}(\{J_{\rho,s}\}_{\rho,s};
\beta;\operatorname{top}(\rho),\operatorname{twp}(s))
\longrightarrow L^k\times([0,1]\times L')
\]
为弱浸没。对 $L^k\times L'$ 因子的弱浸没性
可依第 7.1 节的方法证明。
再添加 $[0,1]$ 方向扩张仓西邻域，
并增加一个一维障碍丛，
就使评价映射的 $[0,1]$ 因子也弱浸没。
于是 (7.2.144) 弱浸没；见注 A1.44。
M. Akaho 与 D. Joyce 在 \textbf{[AkJo08]}
也提到过这一点。

另一类模空间是 (4.6.35) 后定义的
\begin{equation}
\mathcal M_{k+1}^{\mathrm{main},\mathfrak d}
(M',L',\{J_{1,s}\}_s;\beta;
\operatorname{twp}(s);\vec{\mathcal P}),
\tag{7.2.145}\label{eq:7-2-145}
\end{equation}
其中 $\vec{\mathcal P}=(\mathcal P_1,\ldots,
\mathcal P_k)$，
$\mathcal P_i\in\mathcal X_{g_i}([0,1]\times L')$。
置 $\mathfrak d(i)=g_i$，
把 (7.2.145) 简写为
\begin{equation}
\mathcal M_{k+1}^{\mathrm{main},\mathfrak d}
(\{J_{1,s}\}_s;\beta;
\operatorname{twp}(s);\vec{\mathcal P}).
\tag{7.2.146}\label{eq:7-2-146}
\end{equation}
这里省去 $M'$ 与 $L'$。
其评价映射是
\begin{equation}
ev_0^+=(ev_s,ev_0)\colon
\mathcal M_{k+1}^{\mathrm{main},\mathfrak d}
(\{J_{1,s}\}_s;\beta;
\operatorname{twp}(s);\vec{\mathcal P})
\longrightarrow[0,1]\times L'.
\tag{7.2.147}\label{eq:7-2-147}
\end{equation}
见 (4.6.35)。同 (7.2.144) 的弱浸没性一样，
可假设 (7.2.147) 弱浸没。

下面说明这些模空间在 $s=0,1$ 时的行为。
\begin{lemma}{7.2.148}\label{lem:7-2-148}
对 $s_0=0$ 或 $1$，存在映射
\[
\begin{aligned}
\operatorname{Ic}_{s_0}\colon
\mathcal M_{k+1}^{\mathrm{main}}
(\{J_{\rho,s_0}\}_\rho:\beta;
\operatorname{top}(\rho);\vec P)
\\[-4pt]\longrightarrow
\mathcal M_{k+1}^{\mathrm{main}}
(\{J_{\rho,s}\}_{\rho,s};\beta;
\operatorname{top}(\rho),
\operatorname{twp}(s);\vec P),
\end{aligned}
\]
使图式~\ref{diag:7-2-1} 交换。
这里的源模空间（$s_0=0,1$）
用于定义
$\mathfrak f_{s=0}^{\psi_0,\{J^{s=0}_\rho\}_\rho,
\mathfrak s^{\operatorname{top}(\rho)}_{s=0}}$ 与
$\mathfrak f_{s=1}^{\psi_1,\{J^{s=1}_\rho\}_\rho,
\mathfrak s^{\operatorname{top}(\rho)}_{s=1}}$。
$\operatorname{Ic}_{s_0}$ 提升为障碍丛的
逐纤维单射同态，且提升与仓西映射相容。
\end{lemma}

\par\medskip\noindent\begin{minipage}{\textwidth}
\centering
\setlength{\unitlength}{0.75mm}
\begin{picture}(128,55)
\put(19,43){\makebox(0,0){$\mathcal M_{k+1}^{\mathrm{main}}(\{J_{\rho,s_0}\}_\rho:\beta;\operatorname{top}(\rho);\vec P)$}}
\put(105,43){\makebox(0,0){$L'$}}
\put(19,9){\makebox(0,0){$\mathcal M_{k+1}^{\mathrm{main}}(\{J_{\rho,s}\}_{\rho,s};\beta;\operatorname{top}(\rho),\operatorname{twp}(s);\vec P)$}}
\put(105,9){\makebox(0,0){$[0,1]\times L'$}}
\put(72,43){\vector(1,0){26}}
\put(84,46){\makebox(0,0)[b]{$\scriptstyle ev_0$}}
\put(72,9){\vector(1,0){21}}
\put(82,12){\makebox(0,0)[b]{$\scriptstyle ev_0^+$}}
\put(19,36){\vector(0,-1){18}}
\put(24,27){\makebox(0,0)[l]{$\scriptstyle\operatorname{Ic}_{s_0}$}}
\put(105,36){\vector(0,-1){18}}
\put(110,27){\makebox(0,0)[l]{$\scriptstyle x\mapsto(s_0,x)$}}
\end{picture}
\par{\renewcommand{\thefigure}{7.2.1}\refstepcounter{figure}
\normalfont\bfseries 图式~\thefigure\label{diag:7-2-1}}
\end{minipage}\par\medskip

引理~\ref{lem:7-2-148} 由构造直接得出。
\begin{remark}{7.2.149}\label{rem:7-2-149}
引理~\ref{lem:7-2-148} 结论的最后一部分，
也可表述为 $\operatorname{Ic}_{s_0}$
是带仓西结构空间之间的态射。
我们尚未定义这种概念，
一方面是迄今并不需要，
另一方面是其最佳定义对我们仍不明确。
\end{remark}

接着考虑模空间 (7.2.146)。
设 $\vec{\mathcal P}=(\mathcal P_1,\ldots,
\mathcal P_k)$，
$\mathcal P_i\in\mathcal X_{g_i}([0,1]\times L')$。
记
\[
\vec{\mathcal P}|_{s=s_0}
=(\mathcal P_1|_{s=s_0},\ldots,
\mathcal P_k|_{s=s_0}),
\qquad
\mathcal P_i|_{s=s_0}\in
\begin{cases}
\mathcal X^{s=0}_3,&s_0=0,\\
\mathcal X^{s=1}_3,&s_0=1.
\end{cases}
\]
\begin{lemma}{7.2.150}\label{lem:7-2-150}
对 $s_0=0$ 或 $1$，存在映射
\[
\operatorname{Ic}_{s_0}^{\rho=1}\colon
\mathcal M_{k+1}^{\mathrm{main}}
(L',J'_{s_0};\beta;
\vec{\mathcal P}|_{s=s_0})
\longrightarrow
\mathcal M_{k+1}^{\mathrm{main},\mathfrak d}
(\{J_{1,s}\}_s;\beta;
\operatorname{twp}(s);\vec{\mathcal P}),
\]
使图式~\ref{diag:7-2-2} 交换。
它提升为障碍丛的逐纤维单射同态，
且提升与仓西映射相容。
\end{lemma}

\par\medskip\noindent\begin{minipage}{\textwidth}
\centering
\setlength{\unitlength}{0.8mm}
\begin{picture}(125,52)
\put(17,41){\makebox(0,0){$\mathcal M_{k+1}^{\mathrm{main}}(L',J'_{s_0};\beta;\vec{\mathcal P}|_{s=s_0})$}}
\put(101,41){\makebox(0,0){$L'$}}
\put(17,8){\makebox(0,0){$\mathcal M_{k+1}^{\mathrm{main},\mathfrak d}(\{J_{1,s}\}_s;\beta;\operatorname{twp}(s);\vec{\mathcal P})$}}
\put(101,8){\makebox(0,0){$[0,1]\times L'$}}
\put(69,41){\vector(1,0){25}}
\put(81,44){\makebox(0,0)[b]{$\scriptstyle ev_0$}}
\put(69,8){\vector(1,0){20}}
\put(79,11){\makebox(0,0)[b]{$\scriptstyle ev_0^+$}}
\put(17,34){\vector(0,-1){18}}
\put(22,25){\makebox(0,0)[l]{$\scriptstyle\operatorname{Ic}_{s_0}^{\rho=1}$}}
\put(101,34){\vector(0,-1){18}}
\put(106,25){\makebox(0,0)[l]{$\scriptstyle x\mapsto(s_0,x)$}}
\end{picture}
\par{\renewcommand{\thefigure}{7.2.2}\refstepcounter{figure}
\normalfont\bfseries 图式~\thefigure\label{diag:7-2-2}}
\end{minipage}\par\medskip

此处证明同样由构造直接得出。

下面描述所用模空间的边界，先看 (7.2.146)。
设 $\vec{\mathcal P}=(\mathcal P_1,\ldots,
\mathcal P_k)$，
奇异单形 $\mathcal P_i$ 由
\begin{equation}
f_i\colon|\mathcal P_i|\to[0,1]\times L'
\tag{7.2.151}\label{eq:7-2-151}
\end{equation}
实现。写 $f_i=(f_i^s,f_i^{L'})$。
模空间
$\mathcal M_{k+1}^{\mathrm{main},\mathfrak d}
(\{J_{1,s}\}_s;\beta;
\operatorname{twp}(s);\vec{\mathcal P})$
的元素 $\mathbf p$ 可写成
\begin{equation}
\mathbf p=
\bigl(s,(\Sigma,(z_0,\ldots,z_k)),
u,(x_1,\ldots,x_k)\bigr),
\tag{7.2.152.1}\label{eq:7-2-152-1}
\end{equation}
其中 $\Sigma$ 是亏格零带边黎曼曲面，
$(z_0,\ldots,z_k)$ 是按循环序排列的边界标记点，
并满足
\begin{align*}
\textup{(7.2.152.2)}\quad&
u\colon(\Sigma,\partial\Sigma)\to(M',L')
\text{ 是 }J_{1,s}\text{-全纯映射},\\
\textup{(7.2.152.3)}\quad&x_i\in|\mathcal P_i|,\\
\textup{(7.2.152.4)}\quad&f_i^s(x_i)=s
\quad(i=1,\ldots,k),\\
\textup{(7.2.152.5)}\quad&f_i^{L'}(x_i)=u(z_i)
\quad(i=1,\ldots,k).
\end{align*}
回忆评价映射由下式给出：
\begin{equation}
ev_0^+(\mathbf p)=(s,u(z_0)).
\tag{7.2.152.6}\label{eq:7-2-152-6}
\end{equation}
设 $\mathbf p_\alpha$ 是该模空间中趋于边界的点列，
其各分量带下标 $\alpha$。
则有三种可能：
\begin{align*}
\textup{(7.2.153.1)}\quad&
\lim_{\alpha\to\infty}x_{i,\alpha}
\in\partial|\mathcal P_i|,\\
\textup{(7.2.153.2)}\quad&
(\Sigma_\alpha,(z_{0,\alpha},\ldots,z_{k,\alpha}))
\text{ 在 }\alpha\to\infty\text{ 时出现边界冒泡},\\
\textup{(7.2.153.3)}\quad&
s_\alpha\to0\text{ 或 }1
\quad(\alpha\to\infty).
\end{align*}
(7.2.153.2) 可如下描述：
以下说明对应于 $z_{\ell,\alpha},\ldots,z_{m,\alpha}$
落在冒泡上的情形。
取 $\ell,m\in\{1,\ldots,k\}$，
$\ell-1\le m$；其中也允许 $\ell-1=m$。
置
\begin{equation}
Q^0_{\ell,m}=ev_{0*}^+
\bigl(\mathcal M_{m-\ell+2}^{\mathrm{main}}
(\{J_{1,s}\}_s;\beta_1;
\operatorname{twp}(s);
(\mathcal P_\ell,\ldots,\mathcal P_m))\bigr).
\tag{7.2.154.1}\label{eq:7-2-154-1}
\end{equation}
那么，与 (7.2.153.2) 对应的冒泡可表为
\begin{equation}
\mathcal M_{k+\ell-m+1}^{\mathrm{main}}
(\{J_{1,s}\}_s;\beta_2;
\operatorname{twp}(s);
(\mathcal P_1,\ldots,\mathcal P_{\ell-1},
Q^0_{\ell,m},\mathcal P_{m+1},\ldots,
\mathcal P_k)),
\tag{7.2.154.2}\label{eq:7-2-154-2}
\end{equation}
其中 $\beta=\beta_1+\beta_2$；
见图~\ref{fig:7-2-9}。
(7.2.153.3) 中的边界项包含于端点，
故在 $S([0,1]\times L')$ 中为零。

\par\medskip\noindent\begin{minipage}{\textwidth}
\centering
\setlength{\unitlength}{0.8mm}
\begin{picture}(125,67)
\linethickness{0.5pt}
\put(43,38){\circle{35}}
\put(78,38){\circle{35}}
\put(43,38){\makebox(0,0){$\beta_2$}}
\put(78,38){\makebox(0,0){$\beta_1$}}
\put(60.5,38){\circle*{1.5}}
\put(95.5,38){\vector(1,0){11}}
\put(109,37){$\scriptstyle z_0$}
\put(20,58){$\scriptstyle P$}
\put(49,57){$\scriptstyle P_\ell$}
\put(48,13){$\scriptstyle P_m$}
\put(62,57){$\scriptstyle P_{\ell-1}$}
\put(62,13){$\scriptstyle P_{m+1}$}
\put(91,56){$\scriptstyle P_1$}
\put(93,15){$\scriptstyle P_k$}
\put(31,51){\makebox(0,0){$\scriptstyle\times$}}
\put(53,52){\makebox(0,0){$\scriptstyle\times$}}
\put(53,24){\makebox(0,0){$\scriptstyle\times$}}
\put(68,52){\makebox(0,0){$\scriptstyle\times$}}
\put(68,24){\makebox(0,0){$\scriptstyle\times$}}
\put(90,51){\makebox(0,0){$\scriptstyle\times$}}
\put(90,25){\makebox(0,0){$\scriptstyle\times$}}
\put(58,4){\vector(-1,0){20}}
\put(16,1){$\scriptstyle Q^0_{\ell,m}$}
\qbezier[20](27,54)(10,38)(27,17)
\qbezier[20](59,56)(72,38)(59,17)
\qbezier[20](63,56)(51,38)(63,17)
\qbezier[20](96,54)(114,38)(96,17)
\end{picture}
\par{\renewcommand{\thefigure}{7.2.9}\refstepcounter{figure}
\normalfont\bfseries 图~\thefigure\label{fig:7-2-9}}
\end{minipage}\par\medskip

再描述 (7.2.143) 的边界。
设 $\vec P=(P_1,\ldots,P_k)$，
$P_i$ 由 $f_i\colon|P_i|\to L$ 实现。
该模空间的元素可写为
$\widetilde{\mathbf p}=(s,\mathbf p,(x_1,\ldots,x_k))$，
其中
\begin{align*}
\textup{(7.2.155.1)}\quad&s\in[0,1],\\
\textup{(7.2.155.2)}\quad&
\mathbf p\in\mathcal M_{k+1}^{\mathrm{main}}
(M',L',\{J_{\rho,s}\}_\rho:\beta;\operatorname{top}(\rho)),\\
\textup{(7.2.155.3)}\quad&
x_i\in|P_i|,\qquad f_i(x_i)=ev_i(\mathbf p).
\end{align*}
在 (7.2.155.2) 中，对每个 $s\in[0,1]$，
右端都是定义 4.6.8 的单参数
殆复结构族 $\{J_{\rho,s}\}_\rho$ 所给的模空间；
(7.2.155.3) 的 $ev_i$ 在定义 4.6.10 前给出。
评价映射为
\begin{equation}
ev_0^+(s,\mathbf p,(x_1,\ldots,x_k))
=(s,ev_0(\mathbf p))\in[0,1]\times L'.
\tag{7.2.156}\label{eq:7-2-156}
\end{equation}
现在设该模空间中的点列
$(s_\alpha,\mathbf p_\alpha,
(x_{1,\alpha},\ldots,x_{k,\alpha}))$
趋于边界，并写
$\mathbf p_\alpha=
((\Sigma_\alpha^{(i)},\vec z_\alpha^{(i)}),
(u_\alpha^{(i)}),(\rho_\alpha^{(i)}))$；
记号参见定义 4.6.8 与 (4.6.7)。
极限中可能发生以下情况：
\begin{align*}
\textup{(7.2.157.1)}\quad&
\mathbf p_\alpha\text{ 出现 (4.6.9.1) 的退化},\\
\textup{(7.2.157.2)}\quad&
\mathbf p_\alpha\text{ 出现 (4.6.9.2) 的退化},\\
\textup{(7.2.157.3)}\quad&
\mathbf p_\alpha\text{ 出现 (4.6.9.3) 的退化},\\
\textup{(7.2.157.4)}\quad&
\mathbf p_\alpha\text{ 出现 (4.6.9.4) 的退化},\\
\textup{(7.2.157.5)}\quad&
\lim_{\alpha\to\infty}x_{i,\alpha}
\in\partial|P_i|\text{，对某个 }i,\\
\textup{(7.2.157.6)}\quad&
\lim_{\alpha\to\infty}s_\alpha=0\text{ 或 }1.
\end{align*}
同 (4.6.9) 的情形一样，
(7.2.157.1) 对应的边界分量与
(7.2.157.2) 对应的边界分量相消。
(7.2.157.6) 在 $S([0,1]\times L')$ 中为零。
下面描述 (7.2.157.3)、(7.2.157.4) 的边界。

先看 (7.2.157.3)。设
$\rho_\alpha^{(i)}\to0$，
并令 $z_\alpha^{(\ell)},\ldots,
z_\alpha^{(m)}$ 为 $\Sigma_\alpha^{(i)}$
上的标记点。必要时取子列，可假定 $\ell,m$
不随 $\alpha$ 变化；此处 $\ell\le m+1$，
亦包含 $\ell-1=m$，即
$\Sigma_\alpha^{(i)}$ 上没有标记点的情形。
于是该边界分量由
\begin{equation}
\mathcal M_{k-m+\ell+1}^{\mathrm{main}}
(\{J_{\rho,s}\}_{\rho,s};\beta_2;
\operatorname{top}(\rho),\operatorname{twp}(s);
(P_1,\ldots,P_{\ell-1},Q^0_{\ell,m},
P_{m+1},\ldots,P_k))
\tag{7.2.158.1}\label{eq:7-2-158-1}
\end{equation}
描述，其中 $\beta=\beta_1+\beta_2$，并且
\begin{equation}
Q^0_{\ell,m}=ev_{0*}
\bigl(\mathcal M_{m-\ell+2}^{\mathrm{main}}
(L,J;\beta_1;(P_\ell,\ldots,P_m))\bigr).
\tag{7.2.158.2}\label{eq:7-2-158-2}
\end{equation}
接着看 (7.2.157.4)。
若 $\rho_\alpha^{(i)}\to1$，
则由该模空间的定义中的条件 (4.6.7.2)，
$\Sigma_\alpha^{(i)}$ 必须包含第 $0$ 个标记点。
令 $w_1,\ldots,w_j$ 是
$\partial\Sigma_\alpha^{(i)}$ 上的标记点或奇点。
若 $w_j=z_{\ell_j}$，即第 $\ell_j$ 个标记点，
则在 $[0,1]\times L'$ 上定义链
\begin{equation}
\mathcal R^0_{\ell_j,\ell_{j+1}}
=\{\psi_s\}_{s*}(P_{\ell_j}),
\tag{7.2.159.1}\label{eq:7-2-159-1}
\end{equation}
其中 $\{\psi_s\}_{s*}$ 如定义~\ref{def:7-2-140}，
并置 $\beta_{1,j}=\beta_0=0$。
若 $w_j$ 是奇点，考虑以它为根的圆盘树上
全部圆盘与球面的并，
令 $z_{\ell_j},\ldots,z_{\ell_{j+1}-1}$
是该圆盘并上的全部标记点，
$\beta_{1,j}$ 是树上圆盘的同调类之和。
定义
\begin{equation}
\mathcal R^0_{\ell_j,\ell_{j+1}}
=ev_{0*}^+\bigl(
\mathcal M_{\ell_{j+1}-\ell_j+1}^{\mathrm{main}}
(\{J_{\rho,s}\}_{\rho,s};\beta_{1,j};
\operatorname{top}(\rho),\operatorname{twp}(s);
(P_{\ell_j},\ldots,P_{\ell_{j+1}-1}))\bigr).
\tag{7.2.159.2}\label{eq:7-2-159-2}
\end{equation}
于是，(7.2.157.4) 对应的边界分量就是
\translatornote{原书 (7.2.159.3) 将外层模空间的下标
写成带自由指标 $j$ 的 $\ell_{j+1}-\ell_j+1$；
这里有 $m-1$ 条输入链，故依
$\mathcal M_{k+1}$ 约定改为 $m$；见疑误 A-0446。}
\begin{equation}
\mathcal M_m^{\mathrm{main}}
(\{J_{1,s}\}_s;\beta_2;
\operatorname{twp}(s);
(\mathcal R^0_{\ell_1,\ell_2},\ldots,
\mathcal R^0_{\ell_{m-1},\ell_m})).
\tag{7.2.159.3}\label{eq:7-2-159-3}
\end{equation}

现在可以陈述命题 7.2.160。
它是命题~\ref{prop:7-2-35}
与命题~\ref{prop:7-2-100} 的对应版本，
并用于证明命题 4.6.37。
\begin{proposition}{7.2.160}\label{prop:7-2-160}
给定正整数 $\mathcal K$。
存在多值截面系
$\mathfrak s_{\mathfrak d,k,\beta,
\vec{\mathcal P}}^{\operatorname{twp}(s)}$
（$\|(\mathfrak d,\beta)\|\le\mathcal K$）
与
$\mathfrak s_{k,\beta,\vec P}^{
\operatorname{top}(\rho),\operatorname{twp}(s)}$
（$k+\|\beta\|\le\mathcal K$），
以及 $g\le\mathcal K$ 时的
$\mathcal X_g([0,1]\times L')$。
它们具有以下性质。

\noindent\textup{(7.2.161.1)}\quad
$\mathcal X_g([0,1]\times L')$
满足条件~\ref{def:4-6-30}。

\noindent\textup{(7.2.161.2)}\quad
$\mathfrak s_{\mathfrak d,k,\beta,
\vec{\mathcal P}}^{\operatorname{twp}(s)}$
是模空间
\[
\mathcal M_{k+1}^{\mathrm{main},\mathfrak d}
(\{J_{1,s}\}_s;\beta;
\operatorname{twp}(s);\vec{\mathcal P})
\]
的分段光滑多值截面，且横截于零。
$\mathfrak s_{k,\beta,\vec P}^{
\operatorname{top}(\rho),\operatorname{twp}(s)}$
是模空间
\[
\mathcal M_{k+1}^{\mathrm{main}}
(\{J_{\rho,s}\}_{\rho,s};\beta;
\operatorname{top}(\rho),\operatorname{twp}(s);\vec P)
\]
的多值截面，也横截于零。

\noindent\textup{(7.2.161.3)}\quad
若 $\|(\mathfrak d,\beta)\|=g$，则
\[
ev_{0*}^+\bigl(
\mathcal M_{k+1}^{\mathrm{main},\mathfrak d}
(\{J_{1,s}\}_s;\beta;
\operatorname{twp}(s);\vec{\mathcal P})^{
\mathfrak s_{\mathfrak d,k,\beta,
\vec{\mathcal P}}^{\operatorname{twp}(s)}}\bigr)
\]
分解为 $\mathcal X_{(g)}([0,1]\times L')$ 的元素。

\noindent\textup{(7.2.161.4)}\quad
若 $\|\beta\|+k=g$，则
\[
ev_{0*}^+\bigl(
\mathcal M_{k+1}^{\mathrm{main}}
(\{J_{\rho,s}\}_{\rho,s};\beta;
\operatorname{top}(\rho),\operatorname{twp}(s);\vec P)^{
\mathfrak s_{k,\beta,\vec P}^{
\operatorname{top}(\rho),\operatorname{twp}(s)}}\bigr)
\]
分解为 $\mathcal X_{(g)}([0,1]\times L')$ 的元素。

\noindent\textup{(7.2.161.5)}\quad
上述两类多值截面按下文所述意义相容。

\noindent\textup{(7.2.161.6)}\quad
两类多值截面的零集均位于原仓西映射零集的
$\delta$ 邻域内，其中 $\delta>0$，
且当 $\lambda\to\infty$ 时 $\delta\to0$。

\translatornote{原书 (7.2.161.7)--(7.2.161.8)
把两类端点嵌入及对应引理交叉错配；
下文按引理 7.2.148、7.2.150 的实际定义配对，
保留源分项编号；见疑误 A-0447。}
\noindent\textup{(7.2.161.7)}\quad
$\mathfrak s''_{s=0,k,\beta,
\vec{\mathcal P}|_{s=0}}$
（相应地
$\mathfrak s''_{s=1,k,\beta,
\vec{\mathcal P}|_{s=1}}$）
经引理~\ref{lem:7-2-150} 的
$\operatorname{Ic}^{\rho=1}_0$
（相应地 $\operatorname{Ic}^{\rho=1}_1$）
与 $\mathfrak s_{\mathfrak d,k,\beta,
\vec{\mathcal P}}^{\operatorname{twp}(s)}$ 相容。

\noindent\textup{(7.2.161.8)}\quad
$\mathfrak s^{\operatorname{top}(\rho)}_{s=0,k,\beta,\vec P}$
（相应地
$\mathfrak s^{\operatorname{top}(\rho)}_{s=1,k,\beta,\vec P}$）
经引理~\ref{lem:7-2-148} 的
$\operatorname{Ic}_0$
（相应地 $\operatorname{Ic}_1$）
与 $\mathfrak s_{k,\beta,\vec P}^{
\operatorname{top}(\rho),\operatorname{twp}(s)}$ 相容。
\end{proposition}

现在按 (7.2.153)、(7.2.157) 的不同情形，
更明确地说明 (7.2.161.5) 所要求的相容条件。
对 (7.2.153.1)，它是相容条件~\ref{cond:7-2-38}；
对 (7.2.153.2)，由 (7.2.154.1)--(7.2.154.2)，
它是相容条件~\ref{cond:7-2-44}。
对 (7.2.157.3)，由 (7.2.158.1)--(7.2.158.2)，
它是相容条件~\ref{cond:7-2-102}。
对 (7.2.157.4)，由 (7.2.159.1)--(7.2.159.3)，
它是相容条件~\ref{cond:7-2-104}。
\translatornote{原书此处把相容条件 7.2.104
误排为“30.104”；又把下面的“条件 4.6.30”
误称“余维 4.6.30”；见疑误 A-0448。}
对 (7.2.157.5)，仍是相容条件~\ref{cond:7-2-38}。

\begin{proof}
陈述一经明确，证明纲要与命题~\ref{prop:7-2-100}
或命题~\ref{prop:7-2-35} 大体相同。
先按偏序 $<$ 归纳定义所需多值截面与三角剖分。
偏序如定义~\ref{def:7-2-105}；
在该定义中，以
$\mathfrak s_{\mathfrak d,k,\beta,
\vec{\mathcal P}}^{\operatorname{twp}(s)}$
（相应地
$\mathfrak s_{k,\beta,\vec P}^{
\operatorname{top}(\rho),\operatorname{twp}(s)}$）
替换相应的旧多值截面。
按同一方法可证明类似引理 7.2.47、
7.2.107、7.2.120 的结论，
即构造次序是相容的。

接着对 $<$ 归纳。
每一步先由 (7.2.161.7)--(7.2.161.8)
固定 $s=0,1$ 处的多值截面。
注意端点处的横截性已由
命题~\ref{prop:7-2-100} 的 (7.2.101.2) 保证。
再用定理 A1.23（即 \textbf{[FuOn99II]}
的定理 3.11），在保持端点相容性的同时
把横截性延拓到内部。
这里需要对应于引理 7.2.55 的一个引理，
证明方法相同。
因为 $\vec{\mathcal P}$ 中的 $\mathcal P_i$
满足 (4.6.31.3)，
可找到既满足条件 4.6.30、
又延拓边界上既有剖分的光滑三角剖分。
命题~\ref{prop:7-2-160} 得证。
\end{proof}

除 (4.6.38.8) 外，命题 4.6.37
可由命题~\ref{prop:7-2-160} 推出。
令 $C^{\mathcal K}_-([0,1]\times L')$
为 $\mathcal X_{(\mathcal K)}([0,1]\times L')$
自由生成的 $\Lambda_{0,\mathrm{nov}}$ 模的完备化。
依第 4.6.2 小节所述方式应用
命题~\ref{prop:7-2-160}，
并结合命题~\ref{prop:7-2-78}
与命题~\ref{prop:7-2-126} 的证明方法，
\translatornote{原书这里写“命题 7.2.78 与 127”；
本节没有命题 7.2.127，按所引用的运算修正方法改为命题 7.2.126。}
得到 $C^{\mathcal K}_-([0,1]\times L')$
上的滤过 $A_{n''(\lambda),k''(\lambda)}$ 结构
\[
\mathfrak m_{k,\beta}(\vec{\mathcal P})
=ev_{0*}^+\bigl(
\mathcal M_{k+1}^{\mathrm{main},\mathfrak d}
(\{J_{1,s}\}_s;\beta;
\operatorname{twp}(s);\vec{\mathcal P})^{
\mathfrak s_{\mathfrak d,k,\beta,
\vec{\mathcal P}}^{\operatorname{twp}(s)}}\bigr),
\]
以及滤过 $A_{n''(\lambda),k''(\lambda)}$ 同态
\[
\mathfrak f^{\psi_s,\{J_{\rho,s}\}_{\rho,s},
\mathfrak s^{\operatorname{top}(\rho),\operatorname{twp}(s)}}_{k,\beta}
\colon B_kC_1(L;\Lambda_{0,\mathrm{nov}})[1]
\longrightarrow C^{\mathcal K}_-([0,1]\times L'),
\]
其定义为
\[
\begin{aligned}
&\mathfrak f^{\psi_s,\{J_{\rho,s}\}_{\rho,s},
\mathfrak s^{\operatorname{top}(\rho),\operatorname{twp}(s)}}_{k,\beta}
(\vec P)\\
&\qquad=ev_{0*}^+\bigl(
\mathcal M_{k+1}^{\mathrm{main}}
(\{J_{\rho,s}\}_{\rho,s};\beta;
\operatorname{top}(\rho),\operatorname{twp}(s);\vec P)^{
\mathfrak s_{k,\beta,\vec P}^{
\operatorname{top}(\rho),\operatorname{twp}(s)}}\bigr).
\end{aligned}
\]
严格地说，右端只定义
$\mathfrak m^{\mathrm{geo}}_{k,\beta}$
与 $\mathfrak f^{\mathrm{geo}}_{k,\beta}$；
须像命题~\ref{prop:7-2-78}、
\ref{prop:7-2-126} 的证明那样修正，
才得到左端所示映射。
下面定义 $\operatorname{Incl}$，
并在命题中证明 (4.6.38.8)。

\begin{proposition}{7.2.162}\label{prop:7-2-162}
\liftIndex{I261}%
存在包含 $C_-^{\mathcal K}([0,1]\times L')$ 的
$C^{\mathcal K}([0,1]\times L')$，
其上有滤过 $A_{n''(\lambda),k''(\lambda)}$ 结构，
延拓 $C_-^{\mathcal K}([0,1]\times L')$ 上的
对应结构。还存在 $\Lambda_{0,\mathrm{nov}}$
模同态
\[
\operatorname{Incl}_k
=\sum_\beta e^{\mu_{L'}(\beta)/2}
T^{\omega([\beta])}\operatorname{Incl}_{k,\beta}
\colon B_k(C_2(L')[1])
\longrightarrow C^{\mathcal K}([0,1]\times L')[1],
\]
具有以下性质：

\noindent\textup{(7.2.163.1)}\quad
(7.2.135) 定义的
$\operatorname{Incl}_{1,\beta_0}$
与 $\operatorname{Incl}_{k,\beta}$ 一起定义
严格的滤过 $A_{n''(\lambda),k''(\lambda)}$
同态 $\operatorname{Incl}$。

\noindent\textup{(7.2.163.2)}\quad
$\operatorname{Eval}_{s=s_0}\circ
\operatorname{Incl}_{1,\beta_0}=\mathrm{id}$，
而对 $(k,\beta)\ne(1,\beta_0)$，
$\operatorname{Eval}_{s=s_0}\circ
\operatorname{Incl}_{k,\beta}=0$。
\end{proposition}

\par\noindent\textit{命题~\ref{prop:7-2-162} 的证明。}
先描述用于定义 $\operatorname{Incl}_{k,\beta}$
的模空间。考虑常值双参数殆复结构族
$J_{\rho,s}\equiv J$，记为 $\{J\}_{\rho,s}$；
再用不依赖于 $s$ 的映射 $\psi\colon M\to M$，
如前定义
\[
\mathcal M_{k+1}^{\mathrm{main}}
(\{J\}_{\rho,s};\beta;
\operatorname{top}(\rho),\operatorname{twp}(s)).
\]
其评价映射为
\[
(ev^s,ev_0,\ldots,ev_k)\colon
\mathcal M_{k+1}^{\mathrm{main}}
(\{J\}_{\rho,s};\beta;
\operatorname{top}(\rho),\operatorname{twp}(s))
\longrightarrow[0,1]\times(L')^{k+1}.
\]
其中 $ev^s$ 对应参数 $s$。
这是带仓西结构的空间，
其边界由以下两类纤维积及一类端点分量组成：
\begin{equation}
\begin{gathered}
\prod_{i=1}^{\ell}
\mathcal M_{k_i+1}^{\mathrm{main}}
(\{J\}_{\rho,s};\beta_i;
\operatorname{top}(\rho),\operatorname{twp}(s))
\\
\mathbin{\times_{([0,1]\times L')^\ell}}
\mathcal M_{\ell+1}^{\mathrm{main}}
(\{J\}_s;\beta';\operatorname{twp}(s)),
\qquad\beta'+\sum_i\beta_i=\beta.
\end{gathered}
\tag{7.2.164.1}\label{eq:7-2-164-1}
\end{equation}
\begin{equation}
\mathcal M_{k-\ell}^{\mathrm{main}}(J;\beta_2)
\mathbin{\times_{ev_0,L',ev_i}}
\mathcal M_{\ell+1}^{\mathrm{main}}
(\{J\}_{\rho,s};\beta_1;
\operatorname{top}(\rho),\operatorname{twp}(s)),
\tag{7.2.164.2}\label{eq:7-2-164-2}
\end{equation}
\translatornote{原书 (7.2.164.2) 第二个模空间后另有
游离的 $(J:\beta_2)$，与前一因子重复且无句法位置，故未照录。}
其中 $\beta_1+\beta_2=\beta$，
$\ell=1,\ldots,k-1$，$i=1,\ldots,\ell$。
第三类边界分量为
\begin{equation}
\mathcal M_{k+1}^{\mathrm{main}}
(\{J\}_{\rho,s};\beta;
\operatorname{top}(\rho),s=s_0),
\tag{7.2.164.3}\label{eq:7-2-164-3}
\end{equation}
即映射
\[
ev^s\colon
\mathcal M_{k+1}^{\mathrm{main}}
(\{J\}_{\rho,s};\beta;
\operatorname{top}(\rho),\operatorname{twp}(s))
\longrightarrow[0,1]
\]
在 $s_0=0,1$ 处的纤维。
当 $M$ 为一点时，
$\mathcal M_{k+1}^{\mathrm{main}}
(\{J\}_{\rho,s};\beta_0;
\operatorname{top}(\rho),\operatorname{twp}(s))$
就是 $[0,1]\times\mathcal N_{k+1}$；
此时边界分解 (7.2.164.1)--(7.2.164.2)
对应 (7.1.52)，一般情形同理可证。

我们利用该模空间及其对应关系定义
$\operatorname{Incl}_{k,\beta}$。
但为了证明 (7.2.163.2)，
还须适当修改模空间：
对 $(k,\beta)\ne(1,\beta_0)$，
必须使
\[
(\operatorname{Eval}_{s=s_0})_1
\circ\operatorname{Incl}_{k,\beta}=0.
\]
为此如下折叠模空间中 $s=0,1$ 的部分。
忘去 $ev^s$ 的纤维 (7.2.164.3)
内各元素的时间分配，得到映射
\[
\operatorname{Coll}_{s=s_0}\colon
\mathcal M_{k+1}^{\mathrm{main}}
(\{J\}_{\rho,s};\beta;
\operatorname{top}(\rho),s=s_0)
\longrightarrow
\mathcal M_{k+1}^{\mathrm{main}}(J;\beta).
\]
这里 $\{J\}_{\rho,s}$、$\{\psi\}_s$
都是常值族。

\begin{lemma}{7.2.165}\label{lem:7-2-165}
可对仓西邻域取三角剖分，
使 $\operatorname{Coll}_{s=s_0}$ 成为单纯映射。
\end{lemma}
\begin{proof}
按定义域的组合型和各分支映射的同伦类，
把端点模空间分解为
\[
\mathcal M_{k+1}^{\mathrm{main}}
(\{J\}_{\rho,s};\beta;
\operatorname{top}(\rho),s=s_0)
=\bigcup_{\alpha\in A}\operatorname{Int}\mathcal M_\alpha.
\]
每个 $\operatorname{Int}\mathcal M_\alpha$
的闭包 $\mathcal M_\alpha$ 具有带角仓西结构。
对 $\mathbf x\in\operatorname{Int}\mathcal M_\alpha$，
纤维 $\operatorname{Coll}_{s=s_0}^{-1}(\mathbf x)$
是 $\mathbb R^m$ 中的凸多面体
$\Delta_\alpha$，其中 $m$ 是
稳定映射 $\mathbf x$ 的定义域的分支数。
纤维由时间分配 $\rho_i$ 决定，
它们须满足同类关系 $\rho_i\le\rho_j$。
$\Delta_\alpha$ 不依赖于 $\mathbf x$，且
\[
\operatorname{Coll}_{s=s_0}^{-1}
(\operatorname{Int}\mathcal M_\alpha)
=\Delta_\alpha\times
\operatorname{Int}\mathcal M_\alpha.
\]
定义偏序 $<$：$\beta<\alpha$
当且仅当
$\partial\mathcal M_\alpha\supseteq\mathcal M_\beta$。
若 $\beta<\alpha$，
便有 $\Delta_\alpha$ 与 $\Delta_\beta$
某个层的典范等同。
事实上，$\Delta_\alpha$
由 $\Delta_\beta$ 将若干不等式
$\rho_i\le\rho_j$ 改为
$\rho_i=\rho_j$ 得到；
这对应于将 $\mathcal M_\alpha$
中的定义域分支按 $\rho_i,\rho_j$
的时间序胶合成 $\mathcal M_\beta$
中的元素。

对每个 $\beta<\alpha$，置
$\mathcal N_{\alpha,\beta}
=\Delta_\alpha\times\mathcal M_\beta$。
于是有分解
\begin{equation}
\mathcal M_{k+1}^{\mathrm{main}}
(\{J\}_{\rho,s};\beta;
\operatorname{top}(\rho),s=s_0)
=\bigcup_{(\alpha,\beta)}
\mathcal N_{\alpha,\beta}.
\tag{7.2.166}\label{eq:7-2-166}
\end{equation}
在集合 $\{(\alpha,\beta)\mid\beta<\alpha\}$
上定义偏序：$(\alpha,\beta)<(\alpha',\beta')$
当且仅当下列情形之一成立：
\begin{align*}
\textup{(7.2.167.1)}\quad&\beta<\beta',\\
\textup{(7.2.167.2)}\quad&
\beta=\beta',\quad\alpha'<\alpha.
\end{align*}
按此偏序归纳，可对
$\mathcal N_{\alpha,\beta}$ 的仓西邻域
作所需三角剖分，引理得证。
\end{proof}

引理~\ref{lem:7-2-165} 中的剖分可取成
与 (7.2.164) 的分解相容；这由证明得出。
现在考虑模空间
$\mathcal M_{k+1}^{\mathrm{main}}
(\{J\}_{\rho,s};\beta;
\operatorname{top}(\rho),\operatorname{twp}(s))$
与 $\mathcal M_{k+1}^{\mathrm{main}}(J;\beta)$
两个副本的不交并，
再用 $\operatorname{Coll}_{s=0}$、
$\operatorname{Coll}_{s=1}$ 将它们粘合。
所得空间记作
\begin{equation}
\mathcal M_{k+1}^{\mathrm{main},\operatorname{Coll}}
(\{J\}_{\rho,s};\beta;
\operatorname{top}(\rho),\operatorname{twp}(s)).
\tag{7.2.168}\label{eq:7-2-168}
\end{equation}
评价映射 $ev$ 可经 (7.2.168) 分解。
此空间的边界分成两类纤维积
(7.2.169.1)、(7.2.169.2)，
以及 (7.2.169.3) 类型的分量：
\begin{equation}
\begin{gathered}
\prod_{i=1}^{\ell}
\mathcal M_{k_i+1}^{\mathrm{main},\operatorname{Coll}}
(\{J\}_{\rho,s};\beta_i;
\operatorname{top}(\rho),\operatorname{twp}(s))
\\
\mathbin{\times_{([0,1]\times L')^\ell}}
\mathcal M_{\ell+1}^{\mathrm{main},\mathfrak d}
(\{J\}_s;\beta';\operatorname{twp}(s)),
\qquad\beta'+\sum_i\beta_i=\beta.
\end{gathered}
\tag{7.2.169.1}\label{eq:7-2-169-1}
\end{equation}
\begin{equation}
\mathcal M_{k-\ell}^{\mathrm{main}}(J;\beta_2)
\mathbin{\times_{ev_0,L',ev_i}}
\mathcal M_{\ell+1}^{\mathrm{main},\operatorname{Coll}}
(\{J\}_{\rho,s};\beta_1;
\operatorname{top}(\rho),\operatorname{twp}(s)),
\tag{7.2.169.2}\label{eq:7-2-169-2}
\end{equation}
其中 $\beta_1+\beta_2=\beta$，
$\ell=1,\ldots,k-1$，$i=1,\ldots,\ell$；以及
\begin{equation}
\{0,1\}\times\mathcal M_{k+1}^{\mathrm{main}}(J;\beta).
\tag{7.2.169.3}\label{eq:7-2-169-3}
\end{equation}
按 (7.2.168) 的定义，
(7.2.164.3) 变成 (7.2.169.3)。

由于用折叠映射 $\operatorname{Coll}_{s=s_0}$
缩小了 (7.2.168) 中模空间的一部分，
它是否仍有仓西结构并非先验显然。
不过由引理~\ref{lem:7-2-165}，
$\mathcal M_{\ell+1}^{\mathrm{main}}
(\{J\}_{\rho,s};\beta_1;
\operatorname{top}(\rho),\operatorname{twp}(s))$
的仓西邻域可分解成带角流形，
并在
$\mathcal M_{\ell+1}^{\mathrm{main},\operatorname{Coll}}
(\{J\}_{\rho,s};\beta_1;
\operatorname{top}(\rho),\operatorname{twp}(s))$
上诱导相应分解，
使后者的边界由 (7.2.169) 描述。
注意模空间
\[
\mathcal M_{k+1}^{\mathrm{main}}
(\{J\}_{\rho,s};\beta;
\operatorname{top}(\rho),s=s_0)
\]
上的障碍丛与仓西映射是
$\mathcal M_{k+1}^{\mathrm{main}}(J;\beta)$
上相应数据沿 $\operatorname{Coll}_{s=s_0}$ 的拉回；
自同构群（有限群）在
$\operatorname{Coll}_{s=s_0}$ 的纤维上的作用也平凡。
因此可在上述分解的每一分量上
按维数归纳应用定理 A1.23，
仍得到与 (7.2.168) 对应的结论；
见注 7.2.122。

令 $P_i$ 是 $L'$ 上的奇异单形，
$\mathcal P_i$ 是 $[0,1]\times L'$ 上的奇异单形。
与前面一样用纤维积定义
\[
\mathcal M_{k+1}^{\mathrm{main},\operatorname{Coll}}
(\{J\}_{\rho,s};\beta;
\operatorname{top}(\rho),\operatorname{twp}(s);\vec P),
\qquad
\mathcal M_{k+1}^{\mathrm{main},\mathfrak d}
(\{J\}_s;\beta;
\operatorname{twp}(s);\vec{\mathcal P}).
\]

现在陈述证明命题~\ref{prop:7-2-162}
所需的主要技术结果。
令 $\mathcal X_{(\mathcal K)}([0,1]\times L')$
是命题~\ref{prop:7-2-160} 的奇异单形族，
暂写成 $\mathcal X'_{00}([0,1]\times L')$：
也就是暂时忽略命题~\ref{prop:7-2-160}
已分配的代，把所有这些单形看作代零。
命题 7.2.170 将另行分配代。

\begin{proposition}{7.2.170}\label{prop:7-2-170}
给定 $\mathcal K>0$。
存在多值截面系
$\mathfrak s_{\mathfrak d,k,\beta,
\vec{\mathcal P}}^{\operatorname{twp}(s)}$
（$\|(\mathfrak d,\beta)\|\le\mathcal K$）
与
$\mathfrak s_{k,\beta,\vec P,
\operatorname{Incl}}^{
\operatorname{top}(\rho),\operatorname{twp}(s)}$
（$k+\|\beta\|\le\mathcal K$），
及 $g\le\mathcal K$ 时的可数奇异单形集
$\mathcal X'_g([0,1]\times L')$，
满足以下性质。

\noindent\textup{(7.2.171.1)}\quad
$\mathcal X'_g([0,1]\times L')$
满足条件~\ref{def:4-6-30}，
并生成与 $S([0,1]\times L')$
链同伦等价的链复形。

\noindent\textup{(7.2.171.2)}\quad
$\mathfrak s_{\mathfrak d,k,\beta,
\vec{\mathcal P}}^{\operatorname{twp}(s)}$
是模空间
$\mathcal M_{k+1}^{\mathrm{main},\mathfrak d}
(\{J\}_s;\beta;
\operatorname{twp}(s);\vec{\mathcal P})$
横截于零的多值截面。
$\mathfrak s_{k,\beta,\vec P,
\operatorname{Incl}}^{
\operatorname{top}(\rho),\operatorname{twp}(s)}$
是模空间
\[
\mathcal M_{k+1}^{\mathrm{main},\operatorname{Coll}}
(\{J\}_{\rho,s};\beta;
\operatorname{top}(\rho),\operatorname{twp}(s);\vec P)
\]
的分段光滑多值截面，横截于零。

\noindent\textup{(7.2.171.3)}\quad
若 $\|(\mathfrak d,\beta)\|=g$，则
\[
ev_{0*}^+\bigl(
\mathcal M_{k+1}^{\mathrm{main},\mathfrak d}
(\{J\}_s;\beta;
\operatorname{twp}(s);\vec{\mathcal P})^{
\mathfrak s_{\mathfrak d,k,\beta,
\vec{\mathcal P}}^{\operatorname{twp}(s)}}\bigr)
\]
分解为 $\mathcal X'_{(g)}([0,1]\times L')$ 的元素。

\noindent\textup{(7.2.171.4)}\quad
若 $\|\beta\|+k=g$，则
\[
ev_{0*}^+\bigl(
\mathcal M_{k+1}^{\mathrm{main},\operatorname{Coll}}
(\{J\}_{\rho,s};\beta;
\operatorname{top}(\rho),\operatorname{twp}(s);\vec P)^{
\mathfrak s_{k,\beta,\vec P,
\operatorname{Incl}}^{
\operatorname{top}(\rho),\operatorname{twp}(s)}}\bigr)
\]
分解为 $\mathcal X'_{(g)}([0,1]\times L')$ 的元素。

\noindent\textup{(7.2.171.5)}\quad
$\mathfrak s_{\mathfrak d,k,\beta,
\vec{\mathcal P}}^{\operatorname{twp}(s)}$
与
$\mathfrak s_{k,\beta,\vec P,
\operatorname{Incl}}^{
\operatorname{top}(\rho),\operatorname{twp}(s)}$
按下文意义相容。

\noindent\textup{(7.2.171.6)}\quad
二者的零集均位于原仓西映射零集的
$\delta$ 邻域中，其中 $\delta>0$，
且 $\lambda\to\infty$ 时 $\delta\to0$。

\noindent\textup{(7.2.171.7)}\quad
$\mathcal X'_{00}([0,1]\times L')$
包含于 $\mathcal X'_0([0,1]\times L')$；
若 $\mathcal P_i\in\mathcal X'_{00}$，
则 $\mathfrak s_{\mathfrak d,k,\beta,
\vec{\mathcal P}}^{\operatorname{twp}(s)}$
与命题~\ref{prop:7-2-160} 中的截面一致。

\noindent\textup{(7.2.171.8)}\quad
当 $\beta=\beta_0$、$\ell=1$ 时，
还满足下文所述的另一相容条件。

\noindent\textup{(7.2.171.9)}\quad
$\operatorname{Incl}_{1,\beta_0}$ 的像中
出现的全部奇异单形，均属于
$\mathcal X'_0([0,1]\times L')$。
\end{proposition}

下面说明 (7.2.171.5) 的相容条件。
模空间
\[
\mathcal M_{k+1}^{\mathrm{main},\operatorname{Coll}}
(\{J\}_{\rho,s};\beta;
\operatorname{top}(\rho),\operatorname{twp}(s);\vec P)
\]
的边界分成四类。
第一类由某个 $P_i$ 被 $\partial P_i$ 替换而来，
相容条件同相容条件~\ref{cond:7-2-38}。
第二类对应 (7.2.169.1)，
相容条件同相容条件~\ref{cond:7-2-104}。
第三类对应 (7.2.169.2)，
相容条件同相容条件~\ref{cond:7-2-102}。
第四类对应 (7.2.169.3)，
边界分量为
\[
\{0,1\}\times\mathcal M_{k+1}^{\mathrm{main}}
\mathbin{\times_{(L')^{k+1}}}
(P_1\times\cdots\times P_k).
\]
在此要求多值截面及其零集的三角剖分
与定义 $\mathfrak m_{k+1,\beta}(P_1,\ldots,P_k)$
时采用的完全相同。
模空间
$\mathcal M_{k+1}^{\mathrm{main},\mathfrak d}
(\{J\}_s;\beta;
\operatorname{twp}(s);\vec{\mathcal P})$
边界上的相容条件同命题~\ref{prop:7-2-160}。

再说明 (7.2.171.8)。注意
\[
\mathcal M_{1+1}^{\mathrm{main},\operatorname{Coll}}
(\{J\}_{\rho,s};\beta_0;
\operatorname{top}(\rho),\operatorname{twp}(s))
\cong[0,1]\times L'.
\]
这里 $\beta_0=0$；左端就是
$\mathcal N_{1+1}\times([0,1]\times L')$，
而 $\mathcal N_{1+1}$ 为一点。
因此，当 $\beta=\beta_0$、$\ell=1$ 时，
(7.2.169.2) 中的空间
与 $\mathcal M_{k+1}^{\mathrm{main}}(J;\beta)$ 相同，
且
\[
\mathcal M_{1+1}^{\mathrm{main},\operatorname{Coll}}
(\{J\}_{\rho,s};\beta_0;
\operatorname{top}(\rho),\operatorname{twp}(s);P)
\]
横截，并同构于 $[0,1]\times P$。
我们要求该模空间不受扰动，
并采用定义
$\operatorname{Incl}_{1,\beta_0}$
时的相同剖分。

\begin{proof}
命题~\ref{prop:7-2-170} 的陈述一经明确，
证明大部分重复先前同类命题的归纳论证。
归纳第一步是上文的 (7.2.171.8)。
然后定义 $\mathcal X'_0([0,1]\times L')$，
使它既包含 $\mathcal X'_{00}([0,1]\times L')$
又包含 $\operatorname{Incl}_{1,\beta_0}$ 的像，
并满足 (7.2.171.1)。
随后对由 $\mathcal X'_0([0,1]\times L')$
中的奇异单形所取的模空间纤维积
定义多值截面，使 (7.2.171.7) 成立。

注意
$\mathcal M_{k+1}^{\mathrm{main}}
(\{J\}_{\rho,s};\beta;
\operatorname{top}(\rho),s=s_0)$
上的障碍丛与仓西映射，
都是 $\mathcal M_{k+1}^{\mathrm{main}}(J;\beta)$
上对应数据沿 $\operatorname{Coll}_{s=s_0}$ 的拉回；
自同构群（有限群）在折叠映射的纤维上
作用平凡。由此可选取在 $s=s_0$ 处相容的多值截面。
其余证明与命题~\ref{prop:7-2-160} 完全类似，
使用前述同类归纳过程，故略去。
\end{proof}

令 $C^{\mathcal K}([0,1]\times L')$
为 $\mathcal X'_{(\mathcal K)}([0,1]\times L')$
生成的自由 $\Lambda_{0,\mathrm{nov}}$ 模的完备化。
由 (7.2.171.3) 得其滤过
$A_{n''(\lambda),k''(\lambda)}$ 结构，
由 (7.2.171.4) 得滤过
$A_{n''(\lambda),k''(\lambda)}$ 同态
$\operatorname{Incl}\colon
C^{\mathcal K}(L';\Lambda_{0,\mathrm{nov}})
\to C^{\mathcal K}([0,1]\times L')$。
它是同态由 (7.2.171.5) 保证，
而 (7.2.171.8) 给出
$\operatorname{Eval}_{s=s_0}\circ
\operatorname{Incl}=\mathrm{id}$。
更具体地，由于以 $\operatorname{Coll}_{s=s_0}$
折叠模空间，当 $(k,\beta)\ne(1,\beta_0)$ 时，
所有 $\operatorname{Incl}_{k,\beta}$
的贡献在 $s=s_0$ 处均因维数计算而为零。
命题~\ref{prop:7-2-162} 得证。

依第 4.6.2 小节的说明应用
命题~\ref{prop:7-2-162}，便知
$\mathfrak g_{s=0}^{\psi_0,\{J^{s=0}_\rho\}_\rho,
\mathfrak s^{\operatorname{top}(\rho)}_{s=0}}$
与
$\mathfrak g_{s=1}^{\psi_1,\{J^{s=1}_\rho\}_\rho,
\mathfrak s^{\operatorname{top}(\rho)}_{s=1}}$
在 $A_{n''(\lambda),K''(\lambda)}$ 意义下同伦。
这里的两个映射在第 4.6.2 小节开头定义。
定理 4.6.25 的 $A_{n,K}$ 版本得证。

\begin{remark}{7.2.172}\label{rem:7-2-172}
\textup{(1)} 说明为何对
$(k,\beta)\ne(1,\beta_0)$
仍须取非零的 $\operatorname{Incl}_{k,\beta}$。
令 $P_i$ 是 $L'$ 上的单形，考虑纤维积
\begin{equation}
\mathcal M_{k+1}^{\mathrm{main}}
(\{J\}_s;\beta;\operatorname{twp}(s))
\mathbin{\times_{([0,1]\times L')^k}}
\prod_{i=1}^{k}([0,1]\times P_i).
\tag{7.2.173}\label{eq:7-2-173}
\end{equation}
(7.2.173) 有仓西结构。
对每个固定的 $s$，
在此模空间上定义用来给
$C_2(L';\Lambda_{0,\mathrm{nov}})$
构造滤过 $A_\infty$ 结构的多值截面。
因此在 (7.2.173) 上可取沿 $s$ 方向不变、
且横截于零的多值截面。
若能用它定义
$C([0,1]\times L')$ 上的滤过
$A_{n''(\lambda),k''(\lambda)}$ 结构，
那么 $\operatorname{Incl}_{1,\beta_0}$
连同其余分量为零就能构成同态。

但 $[0,1]\times P_i$ 不是单形，
因而不能直接用作奇异链复形
$C([0,1]\times L')$ 的生成元；
必须适当细分，最自然的是棱柱剖分。
剖分之后，能否仍取沿 $s$ 方向不变
并横截于零的多值截面，就不再显然。
故须取在 $s=0,1$ 处相容、
但未必沿 $s$ 不变的多值截面。
若此时只令
$\operatorname{Incl}_{1,\beta_0}$ 非零，
其余 $\operatorname{Incl}_{k,\beta}$ 为零，
便不满足滤过
$A_{n''(\lambda),k''(\lambda)}$ 同态的定义关系。
这迫使我们取非零的高阶分量；
粗略说，它们来自两个多值截面之间的同伦：
一个在棱柱剖分得到的每个单形上零集横截，
另一个沿 $s$ 不变。
因为必须归纳求这些高阶分量，
仍需实行前述过程，
即命题~\ref{prop:7-2-170} 所做的工作。

\textup{(2)} 这是注~\ref{rem:7-2-128} 的续论。
在该注与本小节开头的语境中，
设 $\psi_0=\mathrm{id}$，
$J_\rho^{s=0}\equiv J$ 不依赖于 $\rho$。
此时主张
$\mathfrak g_{s=1}^{\psi_1,\{J_\rho^{s=1}\}_\rho,
\mathfrak s^{\operatorname{top}(\rho)}_{s=1}}$
在 $A_{n,K}$ 意义下同伦于恒等映射。
为证此点，先考虑“顶层”模空间
$\mathcal M_{k+1}^{\mathrm{main}}
(\{J_{\rho,0}\}_\rho:\beta;
\operatorname{top}(\rho);\vec P)$。
因 $J_\rho^{s=0}\equiv J$ 是常值族，
可定义忘却时间分配的映射
\[
\operatorname{forget}\colon
\mathcal M_{k+1}^{\mathrm{main}}
(\{J_{\rho,0}\}_\rho:\beta;
\operatorname{top}(\rho);\vec P)
\longrightarrow
\mathcal M_{k+1}^{\mathrm{main}}(J;\beta;\vec P).
\]
可选取源模空间上的仓西结构与多值截面，
使 $\operatorname{forget}$ 保持它们。
前者的虚维数严格大于后者，
故
$\mathfrak f_{s=0}^{\psi_0,\{J_\rho^{s=0}\}_\rho,
\mathfrak s^{\operatorname{top}(\rho)}_{s=0}}(\vec P)$
是退化单形。然而，要证注~\ref{rem:7-2-128}
的主张，它必须作为奇异链恰等于零。
\translatornote{原书此处将下标 $s=0$ 的 $\mathfrak f$
误配上标 $\psi_1,\{J_\rho^{s=1}\}_\rho,
\mathfrak s_{s=1}^{\operatorname{top}(\rho)}$；
所述顶层模空间及常值族都位于 $s=0$ 端，
故依 (7.2.133) 的定义统一为 $s=0$。}
所以还须修改其定义；
修改方式可仿照命题~\ref{prop:7-2-170} 的证明。

注意，此遗忘映射在扰动后的模空间之间诱导映射
\[
\operatorname{forget}^{\mathfrak s}\colon
\mathcal M_{k+1}^{\mathrm{main}}
(\{J_{\rho,0}\}_\rho:\beta;
\operatorname{top}(\rho);\vec P)^{\mathfrak s}
\longrightarrow
\mathcal M_{k+1}^{\mathrm{main}}
(J;\beta;\vec P)^{\mathfrak s}.
\]
可以选取使 $\operatorname{forget}^{\mathfrak s}$
成为单纯映射的三角剖分。
因此可用该映射将子空间
\[
\mathcal M_{k+1}^{\mathrm{main}}
(\{J_{\rho,0}\}_\rho:\beta;
\operatorname{top}(\rho);\vec P)^{\mathfrak s}
\subset
\mathcal M_{k+1}^{\mathrm{main}}
(\{J_{\rho,s}\}_{\rho,s};\beta;
\operatorname{top}(\rho),\operatorname{twp}(s);\vec P)^{\mathfrak s}
\]
按命题~\ref{prop:7-2-170} 证明中的方式折叠，得到空间
\[
\mathcal M_{k+1}^{\mathrm{main},\operatorname{Coll}}
(\{J_{\rho,s}\}_{\rho,s};\beta;
\operatorname{top}(\rho),\operatorname{twp}(s);\vec P)^{\mathfrak s}.
\]
由于 $\operatorname{forget}^{\mathfrak s}$ 是单纯映射，
这一空间可以三角剖分。
以它替代未折叠的模空间，
便得到从
$\mathfrak g_{s=1}^{\psi_1,\{J^{s=1}_\rho\}_\rho,
\mathfrak s^{\operatorname{top}(\rho)}_{s=1}}$
到恒等映射所需的同伦。
\translatornote{原书此处将“命题 7.2.170”误排为
“Proposition 30.170”；依上下文改正。}

\textup{(3)} 为证明 (3.8.34.1)，
将 (3.8.34.1) 作为其左端的定义，
并在证明 (3.8.33) 的过程中按类似方式
折叠退化链，从而采用这一定义后
(3.8.33) 仍然成立。
\end{remark}

上文定义的滤过 $A_{n,K}$ 代数
$C([0,1]\times L')$
满足 $[0,1]\times C$ 的模型所需的许多公理。
不过其性质与模型公有的性质并不完全相同，
而是略有差别，概括如下。

\begin{definition}{7.2.174}\label{def:7-2-174}
\liftIndex{I264}%
设 $C_{s=s_0}$（$s_0=0,1$）是两个滤过
$A_\infty$ 代数，含有公共子代数 $C_0$；
即 $C_0\subset C_{s=s_0}$，且
$C_{s=s_0}$ 的运算 $\mathfrak m_k$ 限制到
$C_0$ 上恰为 $C_0$ 的 $\mathfrak m_k$。
假设每个嵌入 $C_0\subset C_{s=s_0}$
都是同伦等价。记
$C=(C_0,\{C_{s=s_0}\}_{s_0})$。
若滤过 $A_\infty$ 代数 $\mathfrak C$
连同 $A_\infty$ 同态\liftIndex{I263}\liftIndex{I262}%
\[
\operatorname{Incl}\colon C_0\to\mathfrak C,
\qquad
\operatorname{Eval}_{s=s_0}\colon
\mathfrak C\to C_{s=s_0}
\]
满足以下条件，则称其为 $[0,1]\times C$ 的一个\emph{模型}：
\begin{equation}
\begin{gathered}
(\operatorname{Eval}_{s=s_0})_k=0
\quad(k\ne1),\\
(\operatorname{Eval}_{s=s_0})_1
\text{ 由其 $\mathbb Q$-归约 }
(\overline{\operatorname{Eval}}_{s=s_0})_1
\text{ 诱导},\\
(\operatorname{Eval}_{s=s_0})_{1,\beta}=0
\quad(\beta\ne\beta_0).
\end{gathered}
\tag{7.2.175.1}\label{eq:7-2-175-1}
\end{equation}
\begin{equation}
\operatorname{Eval}_{s=s_0}\circ\operatorname{Incl}
\text{ 是嵌入映射。}
\tag{7.2.175.2}\label{eq:7-2-175-2}
\end{equation}
\begin{equation}
(\operatorname{Eval}_{s=0})_1\oplus
(\operatorname{Eval}_{s=1})_1\colon
\mathfrak C\longrightarrow C_{s=0}\oplus C_{s=1}
\text{ 是满射。}
\tag{7.2.175.3}\label{eq:7-2-175-3}
\end{equation}
\begin{equation}
\overline{\operatorname{Eval}}_{s=s_0}
\text{ 与 }
(\overline{\operatorname{Incl}})_1
\text{ 均为链同伦等价。}
\tag{7.2.175.4}\label{eq:7-2-175-4}
\end{equation}
\end{definition}

取 $C_0=C(L';\Lambda_{0,\mathrm{nov}})$ 与
$C_{s=s_0}=C_3^{s=s_0}(L';\Lambda_{0,\mathrm{nov}})$，
则上文的 $C([0,1]\times L')$ 满足
此定义的 $A_{n,K}$ 版本。

\begin{theorem}{7.2.176}\label{thm:7-2-176}
设 $\mathfrak C^{(i)}$ 是
$[0,1]\times C^{(i)}$ 的模型，
其中 $C^{(i)}=(C_0^{(i)},
\{C_{s=s_0}^{(i)}\}_{s_0})$（$i=1,2$）。
设 $\mathfrak f^{s=s_0}\colon
C_{s=s_0}^{(1)}\to C_{s=s_0}^{(2)}$
是滤过 $A_\infty$ 同态，且
\[
\mathfrak f^{s=0}|_{C_0}
=\mathfrak f^{s=1}|_{C_0}.
\]
假设这些滤过 $A_\infty$ 代数与同态均是有隙的。
则存在滤过 $A_\infty$ 同态
\[
\mathfrak F\colon\mathfrak C^{(1)}\to\mathfrak C^{(2)}
\]
使得
\[
\operatorname{Eval}_{s=s_0}\circ\mathfrak F
=\mathfrak f^{s=s_0}\circ\operatorname{Eval}_{s=s_0},
\qquad
\operatorname{Incl}\circ\mathfrak f^{s=s_0}|_{C_0}
=\mathfrak F\circ\operatorname{Incl}.
\]
\end{theorem}

\begin{proof}
证明大部分与定理 4.2.34 相同，
这里只讨论非滤过版本。
令 $K\to\infty$，归纳构造与
$\operatorname{Eval}$ 和 $\operatorname{Incl}$ 相容的
$A_K$ 同态 $\{\overline{\mathfrak F}_k\}_{k=1}^{K}$。
$\overline{\mathfrak F}_1$ 的构造与第 4.4.1 小节相同。
假设 $k<K$ 时 $\overline{\mathfrak F}_k$ 已构造。
与 $\operatorname{Incl}$ 相容的条件可写成
\[
\overline{\mathfrak F}_K\circ
\overline{\operatorname{Incl}}_1^{\otimes K}
=\text{不含 $\overline{\mathfrak F}_K$ 的项}.
\]
因此可在 $\overline{\operatorname{Incl}}_1^{\otimes K}$
的像上定义 $\overline{\mathfrak F}_K$，
使其与 $\operatorname{Incl}$ 相容；
它在那里也就自动与 $\operatorname{Eval}$ 相容。
现在求 $\overline{\mathfrak F}_K$ 时有一个障碍圈。
它在 $\overline{\operatorname{Incl}}_1^{\otimes K}$
的像上为零，因此是边界，
从而可以找到 $\overline{\mathfrak F}_K$。
证明余部与第 4.4 节相同，故略去。
\end{proof}
